\section{Integral approximation by endpoint features}\label{sec:kernel}

We now approximate the latent matrix \(\mathcal L\) from \eqref{eq:20}
after integrating over its endpoint types.  For an allowed pair
\(k=i+j\) with \(j>0\), let \(S_i,S_k\) be independent site types and
let \(g,h:\Omega_B\to[-1,1]\) be real measurable functions.  The
quantity to be approximated is
\[
 \E_{S_i,S_k}\big[
   \mathcal L_{ik}(S_i,S_k;s)g(S_i)h(S_k)\big].
\]
For a coarse partition supplied by \eqref{eq:32} at a chosen accuracy
\(\epsilon_1\), let \(V_l\) be the features from
Corollary~\ref{cor:channels-features}.  Our target is a finite sum of
scalar multiples of
\[
 \E[g(S)V_l(S)]\,\E[h(S)V_l(S)].
\]
The coefficients may depend on \(B,j,s,l\), but not on the tests.  The
partition is fixed before \(B\) tends to infinity.  The comparison will
be uniform in the tests, allowing a different function at every position.
After summing over allowed pairs and dividing by \(M\), its error will
be controlled by the loss \(e^{-c_3C_*}\) from removing regularity
restrictions, the chosen coarse accuracy \(\epsilon_1\), and a term
tending to zero.  The final proposition records the constants and their
parameter dependence.  Section~\ref{sec:sampling} then converts this integrated
comparison into control for tests chosen after the whole matrix is sampled.

Throughout this section the regularity grid and its tolerances are fixed,
as are \(C_*\), \(\eta>0\), and \(C_6\).
The parameter \(B\) tends to infinity, \(T=\lfloor B^{0.32}\rfloor\),
and \(\eta T\le j\le T\).  All estimates involving the smooth parameter
\(s\) are uniform on the fixed compact interval containing its support.
Constants with a subscript \(\eta\) may depend on \(\eta\); constants
without that subscript in the cutoff-removal estimate below do not.
The parameter \(x\) from the original counting problem does not occur
in this section.  The letters \(x,y\) below denote log coordinates in
\(I=[1/2,3]\).

\subsection{Removing the regularity restrictions}

Let \(S\in\Omega_B\) have the independent-site law: each prime in \(\mathcal P\)
belongs to \(S\) independently with probability \(1/p\).  A fair split
of \(S\) selects each of its primes independently with probability
\(1/2\); write \(b\) for the product of the selected primes.  The
unconditional law of \(b\) is \(\nu_{1/2}\).  For a real measurable
function \(g\) on the site space, define the signed mass
\[
 \gamma_g(n)=\E\big[g(S)\ind_{\{b=n\}}\big].
\]
In particular,
\(\lvert\gamma_g(n)\rvert\le\nu_{1/2}(n)\) whenever
\(\lvert g\rvert\le1\).  All statements below allow \(g\) to depend
on \(B\).

For a subset product \(b\) of \(S\), the untruncated weights satisfy
the exact identity
\[
 A_0(b)K_0(S\setminus b)
 = (\log P_0)R\,2^{-|S|}=B\,2^{-|S|}.
\]
The factor \(2^{-|S|}\) is the conditional probability of each fair
split.  This identity connects the endpoint weights with the unrestricted
splits defining the channels.  The next lemma bounds the cost of removing
the regularity indicators inside the integrated test; after that removal,
the endpoint averages are the unrestricted ones defining the channels.

\begin{lemma}[Uniform removal of cutoffs]\label{lem:kernel-cutoffs}
Let \(k=i+j\) be an allowed pair with \(j>0\), and let
\(\lvert g\rvert,\lvert h\rvert\le1\).  Replacing the coefficient and
remaining-site regularity restrictions in
\(\E[\mathcal L_{ik}g(S_i)h(S_k)]\) by no restrictions, while keeping
the scalar \(k_B\), incurs an absolute error at most
\begin{equation}\label{eq:33}
 C\frac{\Sigma(j)}{T}
       \big(r_B+e^{-c_3C_*}\big),\qquad r_B\longrightarrow0.
\end{equation}
Here \(r_B\) may depend on the fixed regularity parameters and on
\(C_*\), but the constant \(C\) is independent of \(C_*\), \(\eta\),
and \(C_6\).  The pairwise gcd restrictions may then be removed at an
additional error \(O(B^{-998})\).  The resulting integral is
\begin{equation}\label{eq:34}
 k_B B\sum_{a-b=jc}\gamma_g(b)\gamma_h(a)A_0(c)
       \rho_0(\log c/B)
       \Psi_s\big(a/(Tc),b/(Tc)\big).
\end{equation}
The sum is over positive integers \(a,b,c\).
\end{lemma}

\begin{proof}
Applying the subset identity above at both endpoints in \eqref{eq:20}
turns their subset sums into fair-split expectations, still with regularity
indicators attached to the selected and unselected sets, and changes the
coefficient \(k_B/B\) into \(k_BB\).  The remaining coefficient is
\(c=(a-b)/j\), and its weight is \(A_0(c)\) times its regularity
indicator.

All weights other than \(g,h\) are nonnegative.  We may therefore
bound the cost of removing an indicator by putting
\(\lvert g\rvert,\lvert h\rvert\le1\) and summing the corresponding
positive mass.  On the support under consideration,
\(c\asymp X\), \(a,b\asymp TX\), and
\(0.9B\le\log X\le2.2B\).  There are \(O(B)\) dyadic choices of
\(X\).  Formula \eqref{eq:4} bounds the split-product probabilities by
\[
 \nu_{1/2}(a)\nu_{1/2}(b)
 \ll \frac{A_0(a)A_0(b)}{B^2ab}.
\]
Consequently the untruncated mass in one such box is at most
\[
 \frac{C}{B}
 \sum_{\substack{a-b=jc\\c\asymp X,\ a,b\asymp TX}}
       \frac{A_0(a)A_0(b)A_0(c)}{ab}
 \ll \frac{\Sigma(j)}{BT}
\]
by \eqref{eq:7}, since \(ab\asymp T^2X^2\).  If any one of the three
coefficients is nonregular, \eqref{eq:9} gives the same bound with the
additional factor \(r_B+e^{-c_3C_*}\).  Summing over the boxes proves
the required estimate for coefficient failures.

It remains to treat the two unselected endpoint sets.  At a given
prime \(p\), the probabilities of the three possibilities ``selected'',
``unselected'', and ``absent'' are respectively
\(1/(2p),1/(2p),1-1/p\).  Conditional on the selected product being
\(a\), the unselected indicators at primes not dividing \(a\) are
therefore independent with probabilities
\[
 \frac{1/(2p)}{1-1/(2p)}=\frac1{2p-1}.
\]
At primes dividing \(a\) they are zero.  The same statement holds for
\(b\), independently at the other site.  Each selected coefficient on
our support has log size \(O(B)\), and all its prime factors exceed
\(P_0\).  In particular, the reciprocal sum of its prime factors is
\(O(B/P_0)\).  The uniform cutoff probability estimate established
after \eqref{eq:9} thus bounds either unselected-set failure by
\(r_B+O(e^{-c_3C_*})\), uniformly in the chosen coefficients.  Multiplying
by the preceding positive mass bound and summing the boxes proves
\eqref{eq:33}.  This argument uses \eqref{eq:7} and \eqref{eq:9} for
\(1\le j\le T\); it has made no use of \(j\ge\eta T\) or of the block
length \(M\).  This proves the asserted independence of its constant.

Finally, if one of \((a,b),(a,c),(b,c)\) exceeds one, a prime dividing
both \(a\) and \(b\) must occur.  Indeed \(a-b=jc\), and every prime
factor of a coefficient is greater than \(P_0>j\).  The split products
at the two sites are independent, so the probability of a common
selected prime is at most
\[
 \sum_{p\in\mathcal P}\frac1{4p^2}\ll P_0^{-1}.
\]
For a fixed pair of selected products there is at most one value of
\(c\), and the remaining positive integrand is at most
\(C B\sup_c A_0(c)\).  Since
\(\sup A_0\le(\log P_0)R^{1/2}=B^{1/2+o(1)}\), the resulting error is
\(O(B^{3/2+o(1)}/P_0)\).  After removal of these restrictions, the fair
splits and independence of the two sites give \eqref{eq:34} exactly.
\end{proof}

\subsection{Fourier detection and the discarded arcs}

The support of \eqref{eq:34} already predicts the structure of its
real-variable comparison.  For a nonzero term put
\[
 x=\frac{\log a}{B},\qquad y=\frac{\log b}{B},\qquad
 u_1=\frac{a}{Tc},\qquad u_2=\frac{b}{Tc}.
\]
The exact relation \(a-b=jc\) gives
\[
 u_1-u_2=\frac jT,\qquad B(x-y)=\log(u_1/u_2).
\]
Since \(u_1,u_2\) lie in a fixed compact subinterval of \((1,2)\),
the two log coordinates satisfy \(0<x-y\ll B^{-1}\).
We will enforce the discrete relation by additive Fourier orthogonality.
This places the \(c\) sum in a factor to which the arithmetic estimates
apply, while the two signed endpoint sums are controlled by their mass and
\(L^2\) bounds.  For the surviving small denominators, the channel estimates
allow us to average the residue dependence of the endpoint tests with a
uniform error, leaving an explicit finite residue sum.  Fourier inversion
will then give an integral on the narrow logarithmic band, where the coarse
channel approximation produces the feature products described above.

We first analyze \eqref{eq:34} in a single smooth dyadic box.  Choose a
smooth compactly supported dyadic partition of unity on \((0,\infty)\)
for the variable \(c\), and write \(X=2^m\), \(N_X=TX\).  A box has
\(c/X\) in a fixed compact subinterval of \((0,\infty)\).  On the
support of \(\Psi_s\), both \(a/N_X\) and \(b/N_X\) are also in fixed
positive compact intervals.  The smooth weight in the three variables
\[
 (a/N_X,b/N_X,c/X)
\]
has all fixed-order derivatives bounded uniformly in \(m,B,s\).
This includes \(\rho_0((\log X+\log(c/X))/B)\), whose differentiated
log-scale factors only improve the bound.

Extend this weight smoothly inside a larger fixed cube, expand the
extension in a periodic Fourier series, and multiply each coordinate
factor by a fixed smooth compact cutoff equal to one on the original
support.  We obtain a sum of separated products
\[
 w_1(a/N_X)w_2(b/N_X)w_3(c/X).
\]
More precisely, if \(\boldsymbol\nu\in\Z^3\) indexes the Fourier
terms, their scalar coefficients are
\(O_A((1+|\boldsymbol\nu|)^{-A})\) for every fixed \(A\), uniformly
in the boxes and in \(s\).  The fixed-order smooth norms of the factors
grow at most polynomially in \(|\boldsymbol\nu|\).  All estimates below
have only finitely many such smooth-norm losses; choosing \(A\) larger
than those losses plus four makes every Fourier sum absolutely
convergent.  When a bound is summed over boxes, we use the uniform
coefficient bound at each fixed \(\boldsymbol\nu\).  It therefore
suffices to write the calculation for one separated term.

For that term set
\[
 \begin{split}
 S_1(\alpha)&=\sum_a\gamma_h(a)w_1(a/N_X)e(a\alpha),\\
 S_2(\alpha)&=\sum_b\gamma_g(b)w_2(b/N_X)e(b\alpha),\\
 S_3(\alpha)&=\sum_c A_0(c)w_3(c/X)e(c\alpha).
 \end{split}
\]
Additive orthogonality expresses its contribution as
\[
 k_BB\int_0^1 S_1(\alpha)S_2(-\alpha)S_3(-j\alpha)\,d\alpha.
\]
The elementary product formula and Mertens' estimate give
\[
 0\le k_B\le\E K_0(S)
 =R^{1/2}\prod_{p\in\mathcal P}(1-1/(2p))\ll1,
\]
uniformly in the truncation constant.  By \eqref{eq:4},
\eqref{eq:5}, and Parseval,
\begin{equation}\label{eq:35}
 \int_0^1|S_i(\alpha)|^2\,d\alpha
 \ll\frac1{B^2N_X^2}\sum_{n\asymp N_X}A_0(n)^2
 \ll\frac{B^{-2+1/4+o(1)}}{N_X},\qquad i=1,2.
\end{equation}
The first-moment part of \eqref{eq:5} also gives
\begin{equation}\label{eq:kernel-trivial-sums}
 \sup_\alpha|S_i(\alpha)|\ll B^{-1},\qquad i=1,2.
\end{equation}
The implicit constants here and below include the specified smooth
norms of the separated factors.

On the circle of \(j\alpha\pmod1\), take major arcs of radius
\(B^{13}/X\) about the reduced fractions \(u/q\) with
\(q\le B^{12}\).  They are disjoint for large \(B\): the distance
between distinct such fractions is at least \(B^{-24}\), whereas
\(X\) is exponential in \(B\).  On the complement we have
\begin{equation}\label{eq:36}
 |S_3(-j\alpha)|\ll X B^{-1/2+o(1)}.
\end{equation}
Here are the details of the imported estimate and its application.
The Montgomery--Vaughan bound \cite[Corollary~1]{MontgomeryVaughan1977}
states that a multiplicative function \(f\) satisfying
\(\lvert f(p)\rvert\le1\) for every prime \(p\) and
\(\sum_{n\le Z}|f(n)|^2\le Z\) for every \(Z\ge1\) obeys
\[
 \left|\sum_{n\le Y}f(n)e(n\theta)\right|
 \ll \frac{Y}{\log Y}
       +\frac{Y(\log R')^{3/2}}{\sqrt{R'}}
\]
whenever \(\lvert\theta-u'/q'\rvert\le1/(q')^2\),
\((u',q')=1\), and \(2\le R'\le q'\le Y/R'\).
Apply Dirichlet approximation with denominator limit
\(\lfloor X/B^{12}\rfloor\).  If the resulting denominator is at most
\(B^{12}\), the approximation error is at most
\(2B^{12}/X<B^{13}/X\), so the point is on a major arc.  Off the
major arcs we consequently have
\[
 B^{12}<q'\le X/B^{12},\qquad
 |\theta-u'/q'|\le1/(q')^2.
\]
For every \(Y\asymp X\) arising in partial summation, these inequalities
permit \(R'=B^{10}\).  The function
\(f(n)=\mu_{\rm Mob}^2(n)2^{-\omega(n)}\ind_{\{n\ {\rm rough}\}}\)
is multiplicative and 1-bounded, so it satisfies the theorem's
hypotheses.  Smooth partial summation and restoration of the factor
\[
 (\log P_0)R^{1/2}=\sqrt{B\log P_0}=B^{1/2+o(1)}
\]
give \eqref{eq:36}.

For later reference, if a discarded region satisfies
\(\lvert S_3\rvert\ll XB^{-a+o(1)}\), its contribution in one box is
at most
\begin{equation}\label{eq:kernel-discarded-power}
 C B\,XB^{-a+o(1)}
       \frac{B^{-7/4+o(1)}}{TX}
 =\frac{B^{-3/4-a+o(1)}}{T},
\end{equation}
by Cauchy--Schwarz and \eqref{eq:35}.  Summing over the \(O(B)\) boxes
gives a total minor-arc error of
\[
 \frac{B^{-1/4+o(1)}}{T}.
\]

On a major arc, apply \eqref{eq:1} to \(S_3\).  When
\(q>B^{0.4}\), its main term is at most
\(X/\varphi(q)\ll XB^{-0.4+o(1)}\), uniformly along the arc.
Formula \eqref{eq:kernel-discarded-power} shows that these arcs cost
\(B^{-0.15+o(1)}/T\) in total.  The error \(O(XB^{-50})\) in
\eqref{eq:1}, on all the major arcs together, costs at most
\(B^{-49.75+o(1)}/T\).  All three errors are uniform in the endpoint
tests and are \(o(1/T)\).

Put \(Q=B^{0.4}\).  The remaining arcs for \(\alpha\) have the
following exact parametrization:
\[
 d=jq,\qquad h'\pmod d,\qquad (h',q)=1,\qquad
 \alpha=\frac{h'}d+\frac{\xi}{jX},\qquad |\xi|\le B^{13}.
\]
For each reduced fraction on the \(j\alpha\) circle there are exactly
\(j\) lifts in this list.  In particular the parametrization neither
omits nor repeats an arc.  Define
\[
 \widehat w_{3,B,X}(\xi)
 =\int w_3(z)m_B(\log(Xz)/B)e(-\xi z)\,dz.
\]
These transforms are uniformly Schwartz, with bounds controlled by
fixed smooth norms.  The main term from \eqref{eq:1} contains the
factor \(X\), while \(d\alpha=d\xi/(jX)\).  The resulting box
expression is therefore
\begin{equation}\label{eq:37}
 \frac{k_BB}{j}
 \sum_{q\le Q}\frac{\mu_{\rm Mob}(q)}{\varphi(q)}
 \sum_{\substack{h'\bmod d\\(h',q)=1}}
 \int_\R\widehat w_{3,B,X}(\xi)
 S_1\left(\frac{h'}d+\frac{\xi}{jX}\right)
 S_2\left(-\frac{h'}d-\frac{\xi}{jX}\right)\,d\xi.
\end{equation}
In this formula the integration has been extended from
\([-B^{13},B^{13}]\) to \(\R\).  To justify it, use
\eqref{eq:kernel-trivial-sums}, the bound of \(j\varphi(q)\) for the
number of lifts at denominator \(q\), and a Schwartz bound of any
sufficiently large fixed order.  The number of boxes and the total
arc count are polynomial in \(B\), so the tails give \(o(1/T)\) after
all sums.  Notice also that \(d\le TQ\le B^{0.72}<B\), as required
for the channel estimates.

\subsection{Fine histograms and residue averaging}

The discarded arcs already have negligible total contribution. The
remaining major arcs carry a logarithmic coordinate and a unit residue at
each endpoint. We next use the channel bounds to retain the former while
averaging the latter, uniformly for the endpoint tests.

We replace the split-product measures in \eqref{eq:37} by the
fine log and residue histograms of \eqref{eq:30}.  For instance, the
replacement for its first sum is
\begin{equation}\label{eq:38}
 \frac1{\varphi(d)}\sum_{r\in\mathcal R_d}e(h'r/d)
 \int_I U_dh(x,r)w_1(e^{Bx}/N_X)
                   e(\xi e^{Bx}/(jX))\,dx.
\end{equation}
The residue phases have not been approximated.  We give the norm
estimate which controls this step, its denominator sum, and the
subsequent residue projection.

Let \(I_X\subset I\) be an enlarged log window about
\(\log N_X/B\), of length \(O(1/B)\), which contains all fine cells
meeting the support of the factors in a box.  These windows have
bounded overlap as \(X\) ranges over its dyadic values, because
successive centers are spaced by \((\log2)/B\).  They lie in the
interior of \(I\) for large \(B\).  If \(F\) is supported in
\(I_X\) and bounded, finite Fourier Parseval on \(\Z/d\Z\), with
the residue vector extended by zero off the units, gives
\begin{equation}\label{eq:39}
 \begin{split}
 &\sum_{h'\bmod d}\left|
  \frac1{\varphi(d)}\sum_{r\in\mathcal R_d}e(h'r/d)
                    \int_{I_X}V(x,r)F(x)\,dx\right|^2\\
 &\hspace{12mm}\le
  \frac d{\varphi(d)}\,|I_X|\,\|F\|_\infty^2
             \|V\|_{L^2(I_X\times\mathcal R_d)}^2.
 \end{split}
\end{equation}
Indeed the exact finite Fourier identity, before applying
Cauchy--Schwarz to the integrals, is
\[
 \sum_{h'\bmod d}
 \left|\frac1{\varphi(d)}\sum_{r\in\mathcal R_d}e(h'r/d)v_r\right|^2
 =\frac d{\varphi(d)}\frac1{\varphi(d)}
                    \sum_{r\in\mathcal R_d}|v_r|^2.
\]
Here and throughout, the residue component of the \(L^2\) norm uses
uniform probability on \(\mathcal R_d\).

For clarity, the histogram replacement can be estimated even when
the signed masses \(\gamma_g\) have no regularity whatsoever.
For a fine cell \(J\) and a unit residue \(r\), the definition
\eqref{eq:30} says exactly that
\[
 \frac1{\varphi(d)}\int_J U_dh(x,r)\,dx
 =\E\big[h(S)\ind_{\{\log b/B\in J,\ b\equiv r\ (d)\}}\big].
\]
For
\(F(x)=w_1(e^{Bx}/N_X)e(\xi e^{Bx}/(jX))\), differentiation on the
relevant enlarged window gives
\[
 \|F'\|_\infty\le C_{w,\eta}(1+|\xi|)B.
\]
We used \(e^{Bx}\asymp TX\) and \(j\ge\eta T\).  Thus the
oscillation of \(F\) on a fine cell is at most
\[
 a_B(\xi):=C_{w,\eta}(1+|\xi|)B\delta_B.
\]
After taking the factor \(1/\varphi(d)\) outside as in
\eqref{eq:38}, the residue vector of the replacement error has
absolute value at residue \(r\) at most
\[
 a_B(\xi)\int_{I_X}U_d1(x,r)\,dx.
\]
This follows by subtracting the cell average of \(F\) from its value
at the actual split product and using
\(\lvert h\rvert\le1\).  Positivity of the measure for \(U_d1\)
is the only pointwise information used.  Formula \eqref{eq:39}
therefore bounds the squared sum of this error over \(h'\) by
\[
 C\frac d{\varphi(d)}|I_X|\,a_B(\xi)^2
          \|U_d1\|_{L^2(I_X\times\mathcal R_d)}^2.
\]
The actual, unreplaced sum has the same type of estimate without
\(a_B\), by bounding its residue vector with
\(\|F\|_\infty\int_{I_X}U_d1\).  The replaced sum has
\eqref{eq:39} directly with \(V=U_dh\).

We apply these estimates to the difference of the two products in
\eqref{eq:37}, replacing one factor at a time.  Cauchy--Schwarz over
\(h'\) is valid even though the sum is restricted by \((h',q)=1\),
since extending either squared sum to all residues increases it.
Afterwards Cauchy--Schwarz over the boxes and bounded overlap of
\(I_X\) bound the sum of products of local norms by the product of
global norms.  These global norms are bounded by \eqref{eq:31}.
Finally the Schwartz moments absorb the factors \(1+|\xi|\), as
well as the polynomial smooth losses in the separated expansion.
Since \(|I_X|\ll1/B\), this proves that at a fixed \(q\) the total
histogram error, including all boxes, is at most
\begin{equation}\label{eq:kernel-histogram-q}
 \frac{C_\eta}{j}
    \frac1{\varphi(q)}\frac{jq}{\varphi(jq)}\,B\delta_B.
\end{equation}
In particular, the number \(O(B)\) of boxes has not introduced an
extra factor of \(B\): the local window length cancels the \(B\) in
\eqref{eq:37}, and the remaining local norms are summed by bounded
overlap.

We record explicitly a bound for the denominator sum.  The elementary
inequality
\[
 \frac{jq}{\varphi(jq)}
 \le\frac j{\varphi(j)}\frac q{\varphi(q)}
\]
and partial summation imply
\begin{equation}\label{eq:kernel-denominator-sum}
 \sum_{q\le Q}\frac1{\varphi(q)}\frac{jq}{\varphi(jq)}
 \ll\frac j{\varphi(j)}\log(2Q).
\end{equation}
For completeness, to verify the required mean-value input write
\((n/\varphi(n))^2=\sum_{d\mid n}b(d)\), where \(b\) is supported
on squarefree integers and
\[
 b(p)=(1-1/p)^{-2}-1=O(1/p).
\]
All these coefficients are nonnegative, and
\[
 \sum_d\frac{b(d)}d=\prod_p(1+b(p)/p)<\infty.
\]
Hence
\(\sum_{q\le Q}(q/\varphi(q))^2\ll Q\), which after partial
summation gives \eqref{eq:kernel-denominator-sum}.  Since
\(j/\varphi(j)=B^{o(1)}\) uniformly for \(j\le T\) and
\(B\delta_B\asymp B^{-1/10}\), summing
\eqref{eq:kernel-histogram-q} gives \(o(1/T)\), uniformly on the
allowed lag range.

We now replace \(U_dh,U_dg\) in the histogram expression by their
residue averages \(u_h,u_g\).  Expand the difference of products
one factor at a time and repeat \eqref{eq:39}.  The only new input
is
\[
 \|U_dh-u_h\|_2+\|U_dg-u_g\|_2\ll B^{-c_7}
\]
from \eqref{eq:31}, uniformly in \(d\le B\) and the bounded tests.
It follows that the error is bounded by
\[
 \frac{C_\eta}{j}\frac j{\varphi(j)}
                  \log(2Q)B^{-c_7}=o(1/T).
\]
This establishes the residue projection for arbitrary bounded
single-site tests, rather than just for the constant test.

\subsection{The exact residue factor}

After the channel projection, the endpoint tests no longer depend on a
residue coordinate. The remaining finite residue sum can therefore be
computed exactly; it is the arithmetic factor governing the allowed lags.

After residue averaging, the unit phase average in
\eqref{eq:38} is \(c_d(h')/\varphi(d)\), where
\[
 c_d(h')=\sum_{r\in\mathcal R_d}e(h'r/d)
\]
is the Ramanujan sum.  The second endpoint contributes its complex
conjugate phase average.  Since the coefficient of a nonsquarefree
\(q\) in \eqref{eq:37} is zero, only squarefree \(q\) matter.
For such \(q\) the exact identity is
\begin{equation}\label{eq:40}
 \sum_{\substack{h'\bmod jq\\(h',q)=1}}
       \frac{|c_{jq}(h')|^2}{\varphi(jq)^2}
 =\begin{cases}
 \displaystyle\frac{j}{\varphi(j)\varphi(q)},&(j,q)=1,\\[4pt]
 0,&(j,q)>1.
 \end{cases}
\end{equation}
Here is a direct proof that also covers nonsquarefree \(j\).  If
\(p\mid(j,q)\), the modulus \(jq\) is divisible by \(p^2\), while
\((h',q)=1\) forces \(p\nmid h'\).  For every \(a\ge2\),
\[
 c_{p^a}(h')
 =\sum_{r\bmod p^a}e(h'r/p^a)
   -\sum_{r\bmod p^{a-1}}e(h'r/p^{a-1})=0
 \quad(p\nmid h').
\]
Chinese remaindering therefore makes the whole Ramanujan sum zero.
If \((j,q)=1\), the Ramanujan sums split over the two coprime
moduli.  At the squarefree modulus \(q\), a unit frequency has
\(c_q(h')=\mu_{\rm Mob}(q)\), of absolute value one.  At modulus
\(j\), orthogonality gives
\[
 \sum_{h'\bmod j}|c_j(h')|^2
 =\sum_{r,r'\in\mathcal R_j}
        \sum_{h'\bmod j}e(h'(r-r')/j)=j\varphi(j).
\]
The numerator on the left of \eqref{eq:40} is thus
\(j\varphi(j)\varphi(q)\), and division by
\(\varphi(jq)^2=\varphi(j)^2\varphi(q)^2\) proves the identity.

Combining \eqref{eq:40} with the factor
\(\mu_{\rm Mob}(q)/\varphi(q)\) in \eqref{eq:37} yields the
singular series
\begin{equation}\label{eq:41}
 \mathfrak S(j)
 =\frac j{\varphi(j)}
       \sum_{(q,j)=1}\frac{\mu_{\rm Mob}(q)}{\varphi(q)^2}
 =\frac j{\varphi(j)}
       \prod_{p\nmid j}\left(1-\frac1{(p-1)^2}\right).
\end{equation}
The series is absolutely convergent: for every fixed
\(\epsilon>0\), \(1/\varphi(q)\ll_\epsilon q^{-1+\epsilon}\), so
\[
 \sum_{q>Q}\frac1{\varphi(q)^2}
 \ll_\epsilon Q^{-1+2\epsilon}
 \quad(0<\epsilon<1/2).
\]
This tail bound is independent of \(j\).  Also
\[
 0\le\mathfrak S(j)\le j/\varphi(j).
\]
The series vanishes when \(j\) is odd, consistently with the parity
obstruction in \(a-b=jc\) for rough coefficients.

After the residue projection, the archimedean integrals in
\eqref{eq:37} no longer depend on \(q\).  The total over boxes of
their absolute values, before the factor \(B/j\), is \(O(1/B)\):
apply Cauchy--Schwarz on their log windows, then use bounded overlap
and \(\|u_g\|_2,\|u_h\|_2\ll1\).  The Schwartz integrations and
smooth expansions have the same summability as before.  It follows
that replacing the sum \(q\le Q\) by \eqref{eq:41} costs at most
\begin{equation}\label{eq:kernel-series-tail}
 \frac Cj\frac j{\varphi(j)}
                  \sum_{q>Q}\frac1{\varphi(q)^2}.
\end{equation}
In particular this is a uniform vanishing multiple of
\((j/\varphi(j))/j\).

\subsection{Archimedean inversion and coarse features}

For one separated box, the phase remaining in the product of
archimedean integrals is
\[
 e\left(\xi\frac{e^{Bx}-e^{By}}{jX}\right).
\]
Fourier inversion in \(\xi\) therefore evaluates
\(w_3(z')m_B(\log(Xz')/B)\) at
\(z'=(e^{Bx}-e^{By})/(jX)\).  The inversion introduces no further
Jacobian: its integration variable is \(\xi\).  Recombining the
separated factors and the dyadic partition gives
\begin{equation}\label{eq:42}
 \begin{split}
 &\frac{k_BB}{j}\mathfrak S(j)
  \int_{\substack{x>y\\x,y\in I}}u_h(x)u_g(y)m_B(z)\rho_0(z)\\
 &\hspace{14mm}\cdot
  \Psi_s\left(
       \frac{\lambda e^{B(x-y)}}{e^{B(x-y)}-1},
       \frac{\lambda}{e^{B(x-y)}-1}\right)\,dx\,dy,\\
 &\hspace{8mm}\lambda=\frac jT,\qquad
 z=\frac1B\log\left(\frac{e^{Bx}-e^{By}}j\right).
 \end{split}
\end{equation}
Weights are interpreted as zero off their positive domains and
supports.  The inversions and rearrangements are justified for
\(u_g,u_h\in L^2(I)\) by the bounded log windows, the Schwartz
transforms, and the absolutely summable smooth expansion, with the
absolute-product bound just used for \eqref{eq:kernel-series-tail}.

We spell out the support properties needed to approximate
\eqref{eq:42}.  There is a compact interval \([a_0,b_0]\subset(1,2)\)
containing the support of \(\psi\).  Write \(v=B(x-y)\).  If the
\(\Psi_s\) factor is nonzero, its two arguments \(u_1,u_2\) lie
in \([a_0,b_0]\), with
\[
 u_1-u_2=\lambda,\qquad e^v=u_1/u_2.
\]
For \(\eta\le\lambda\le1\), this implies
\begin{equation}\label{eq:kernel-v-support}
 0<c_\eta\le v\le C,
 \qquad c_\eta=\log(1+\eta/b_0),\quad
 C=\log(b_0/a_0).
\end{equation}
If these inequalities have no solution for a particular \(\lambda\),
the weight is simply zero.  On this support,
\[
 z=y+\frac{\log(e^v-1)-\log j}{B}
   =y+O_\eta(\log T/B).
\]
The functions \(m_B\rho_0\), extended by zero beyond a fixed compact
interval inside \((0,4)\), converge uniformly to \(m\rho_0\), with
uniformly bounded derivatives.  It follows that
\[
 m_B(z)\rho_0(z)=m(y)\rho_0(y)+o_\eta(1)
\]
uniformly wherever the integrand can be nonzero.

To quantify its effect for arbitrary tests, consider the integral
operator with nonnegative kernel
\[
 B\ind_{\{0<x-y<C/B\}},\qquad x,y\in I.
\]
Its row and column integrals are at most \(C\), so Schur's inequality
gives a uniform \(L^2\) operator bound.  Multiplication by the bounded
smooth weights preserves that bound.  Since
\(\|u_g\|_2,\|u_h\|_2\ll1\), the replacement of \(m_B(z)\rho_0(z)\)
in \eqref{eq:42} costs
\(o_\eta(1)\mathfrak S(j)/j\).

Given \(\epsilon_1>0\), choose the fixed coarse partition in
\eqref{eq:32}, and write
\(u_g^{\rm c}=P_{m_1}u_g\), \(u_h^{\rm c}=P_{m_1}u_h\).
The same operator bound and the contraction property of
\(P_{m_1}\) show that replacing the two endpoint functions by
\(u_g^{\rm c},u_h^{\rm c}\) costs at most
\[
 C_\eta(\epsilon_1+o(1))\frac{\mathfrak S(j)}j.
\]
This constant does not grow with the chosen coarse partition;
the errors tending to zero may of course depend on that fixed
partition.  Its coarse functions satisfy
\[
 \|u_g^{\rm c}\|_\infty,\ \|u_h^{\rm c}\|_\infty
 \le \max_l |I_l|^{-1/2}
       \max(\|u_g\|_2,\|u_h\|_2)\ll_{m_1}1.
\]
By \eqref{eq:kernel-v-support}, \(x\) and \(y\) belong to the same
coarse cell except when \(y\) is within \(C/B\) of one of the
finitely many cell boundaries.  The area of these exceptional
pairs is \(O_{m_1}(B^{-2})\): the total possible length for \(y\)
is \(O_{m_1}(B^{-1})\), and for each such \(y\) the possible length
for \(x\) is \(O(B^{-1})\).  Thus replacing
\(u_h^{\rm c}(x)\) by \(u_h^{\rm c}(y)\) costs
\(O_{\eta,m_1}(B^{-1})\mathfrak S(j)/j\).  The support of
\(m(y)\rho_0(y)\) is a fixed positive distance from the boundary of
\(I\), so there is no further boundary contribution from the
condition \(x\in I\) for large \(B\).

Recall the coarse features from Corollary~\ref{cor:channels-features},
and define their deterministic coefficients by
\[
 V_l(S)=\frac1{|I_l|}
          \Prob_{\rm split}(\log b/B\in I_l\mid S),
 \qquad h_l=\int_{I_l}m(y)\rho_0(y)\,dy.
\]
These are nonnegative, \(V_l(S)\le |I_l|^{-1}\), and \(h_l\ge0\).
Since the fine partition refines the coarse partition, averaging
\eqref{eq:30} over a coarse cell gives the exact identity
\[
 (P_{m_1}u_g)|_{I_l}=\E[g(S)V_l(S)].
\]
After the preceding replacements in \eqref{eq:42}, substitute
\(x=y+v/B\).  The factor \(dx=dv/B\) cancels its prefactor \(B\),
and \(1/j=(1/T)(1/\lambda)\).  The resulting integral tests are
therefore precisely those of the following site matrix:
\begin{equation}\label{eq:43}
 \begin{split}
 \mathcal M_{ik}(S_i,S_k;s)
 &=\frac{k_B}{T}\mathfrak S(|j|)
       w(s,|j|/T)\sum_l h_lV_l(S_i)V_l(S_k),\qquad j=k-i,\\
 w(s,\lambda)
 &=\frac1\lambda\int_{v>0}
   \Psi_s\left(\frac{\lambda e^v}{e^v-1},
               \frac{\lambda}{e^v-1}\right)\,dv,
       \qquad \eta\le\lambda\le1.
 \end{split}
\end{equation}
The matrix is zero off the allowed pairs.  Formula \eqref{eq:43}
defines it for both signs of \(j\); it is symmetric because its
displayed dependence on the endpoints is symmetric.  For positive
\(j\), independence of the endpoint types turns its integral test
against \(g,h\) into the product of the two expectations just
displayed.  The case of negative \(j\) follows by exchanging the
endpoints.

The function \(w\) is nonnegative and smooth, with all fixed-order
derivatives bounded on the relevant compact \(s\)-range and
\(\eta\le\lambda\le1\).  There is in fact an expression giving
bounds independent of \(\eta\).  In \eqref{eq:43} put
\(u=\lambda/(e^v-1)\).  Then
\(dv=-\lambda\,du/(u(u+\lambda))\), and hence
\begin{equation}\label{eq:kernel-w-ratio}
 w(s,\lambda)=\int_0^\infty
           \frac{\Psi_s(u+\lambda,u)}{u(u+\lambda)}\,du.
\end{equation}
The factors \(\psi(u)\psi(u+\lambda)\) restrict the integration to
a fixed compact subset of \(u>0\).  Consequently this formula
extends \(w\) smoothly to \(0\le\lambda\le1\), with uniformly
bounded derivatives there, and gives zero when
\(\lambda>b_0-a_0\).  Differentiation under the integral is valid
because its denominators are bounded away from zero on that fixed
support.

\subsection{Uniform integral cut comparison and row bounds}

The preceding calculation has produced the finite-feature matrix
$\mathcal M$. We now collect its approximation error in the same block
normalization as the original energy and record the row bounds needed
when the endpoint types are sampled.

\begin{proposition}[Integral comparison]\label{prop:kernel-comparison}
For the matrix \(\mathcal M\) in \eqref{eq:43}, put
\(\mathcal E_{ik}=\mathcal L_{ik}-\mathcal M_{ik}\), with all three
matrices zero on nonpairs.  Given \(\epsilon_1>0\), choose the coarse
partition as in \eqref{eq:32}.  Then
\begin{equation}\label{eq:44}
 \sup_{\substack{g_i,h_k:\Omega_B\to[-1,1]\\
                  g_i,h_k\ \text{real measurable}}}
 \frac1M\left|
  \sum_{i,k}\E_{S_i,S_k}
      [\mathcal E_{ik}(S_i,S_k;s)g_i(S_i)h_k(S_k)]\right|
 \le C e^{-c_3C_*}+C_\eta\epsilon_1+o(1).
\end{equation}
The \(o(1)\) is uniform in \(s\) and in all the indicated tests.
The coefficient of \(e^{-c_3C_*}\) is independent of \(\eta\),
\(C_6\), and \(C_*\).  In addition, for every realization of the
site types,
\begin{equation}\label{eq:kernel-main-row-bounds}
 \sup_i\sum_k|\mathcal M_{ik}|\ll_{\eta,m_1}1,
 \qquad
 \sup_i\sum_k|\mathcal M_{ik}|^2\ll_{\eta,m_1}T^{-1}.
\end{equation}
\end{proposition}

\begin{proof}
Every estimate above was uniform in a pair of tests with absolute
value at most one.  Apart from \eqref{eq:33}, the pre-projection
arc errors are uniform \(o(1/T)\).  The histogram and residue
projection errors have the explicit common upper bound
\[
 \frac{C_\eta}{j}\frac j{\varphi(j)}\log(2Q)
                   (B\delta_B+B^{-c_7}).
\]
The series tail is bounded by \eqref{eq:kernel-series-tail}, and
the archimedean and coarse replacements have total error
\[
 \left(C_\eta\epsilon_1+o_{\eta,m_1}(1)\right)
                  \frac{\mathfrak S(j)}j.
\]
All these bounds hold separately for every allowed edge.  They
thus remain valid if the endpoint tests differ from edge to edge,
and in particular for the collection of tests in \eqref{eq:44}.

For a fixed positive lag there are at most \(M\) ordered pairs
with that lag; including the negative lag at most doubles this
number.  Division by \(M\) therefore reduces their total error
to at most twice the sum of the edge error over
\(\eta T\le j\le T\).  The first-moment bound for \(\Sigma\) gives
\[
 \frac1T\sum_{1\le j\le T}\Sigma(j)\ll1.
\]
Consequently the sum of \eqref{eq:33} is
\(O(r_B+e^{-c_3C_*})\), with a constant independent of
\(\eta,C_6,C_*\).  A uniform \(o(1/T)\) error sums to \(o(1)\).
For the other errors use the bounded mean of \(j/\varphi(j)\)
and
\[
 \sum_{\eta T\le j\le T}\frac1j\frac j{\varphi(j)}
 \le\frac1{\eta T}\sum_{j\le T}\frac j{\varphi(j)}
 \ll_\eta1.
\]
Since \(\log(2Q)(B\delta_B+B^{-c_7})\to0\), the denominator
sum contributes \(o(1)\).  Absolute convergence makes the series
tail tend to zero as well.  Finally
\(\mathfrak S(j)\le j/\varphi(j)\) controls the archimedean and
coarse errors, giving the stated \(C_\eta\epsilon_1+o(1)\).
There is no factor of \(C_6\), because the count of pairs for
each lag has already been divided by \(M\).  This proves
\eqref{eq:44}.

For the deterministic row estimates, boundedness of \(k_B\),
\(w\), and the finitely many features gives the pointwise bound
\[
 |\mathcal M_{ik}|
 \le\frac{C_{\eta,m_1}}T\frac{|k-i|}{\varphi(|k-i|)}
\]
on allowed pairs.  Each row has at most two entries for each
positive lag.  The first and second powers of \(j/\varphi(j)\)
have bounded means.  For the second power this was proved before
\eqref{eq:kernel-denominator-sum}, and the first follows by
Cauchy--Schwarz.  Summing the displayed pointwise bound and its
square over \(1\le j\le T\) proves
\eqref{eq:kernel-main-row-bounds}, uniformly in every collection
of site types and in the block length.
\end{proof}
